\begin{table*}[h]
\centering
\renewcommand{\arraystretch}{1.2}
\setlength{\tabcolsep}{10pt}
\footnotesize
\caption{\textbf{Acoustic Features.}
Vowel-level prosodic features extracted from time-aligned vowel segments, serve as sufficient statistics for speech naturalness judgment.}
\vspace{-0.6em}

\begin{tabular}{p{2.2cm} p{12.6cm}}
\toprule
\textbf{Acoustic Feature} & \textbf{Definition and Interpretation} \\
\midrule
\textbf{Pitch Level} &
Average fundamental frequency ($F_0$) of a vowel segment, reflecting overall pitch height. \\

\textbf{Pitch Var.} &
Within-segment variability of frequency, reflecting the degree of expressiveness. \\

\textbf{Pitch Slope} &
Trend of frequency change across the vowel segment (rising or falling), reflecting local intonation contours. \\

\textbf{Intensity Level} &
Average loudness (e.g., RMS energy) of the vowel segment, reflecting perceived vocal strength or emphasis. \\

\textbf{Intensity Var.} &
Temporal fluctuation of intensity within the vowel segment, indicating stability or variability in loudness. \\

\textbf{Duration} &
Temporal length of the vowel segment, reflecting speech rate, emphasis, and pacing characteristics. \\

\bottomrule
\end{tabular}
\label{table:acoustic-features}
\vspace{-0.5em}
\end{table*}


\begin{table*}[!t]
\centering
\renewcommand{\arraystretch}{1.25}
\setlength{\tabcolsep}{10pt}
\captionsetup{font=small}
\footnotesize
\caption{\textbf{Evidence Log Dimensions.}
Each utterance is analyzed along four dimensions to synthesize interpretable reasoning grounded in acoustic features and transcript context.}
\vspace{-0.6em}

\begin{tabular}{p{2.4cm} p{12.6cm}}
\toprule
\textbf{Dimension} & \textbf{Description} \\
\midrule
\textbf{Inferred Context} &
A concise description of the communicative intent of the utterance (e.g., question, acknowledgment, neutral statement), inferred from the transcript. \\

\textbf{NSVs / Fillers} &
Identification of any non-speech vocalizations or filler words (e.g., ``yeah'', ``uh'',  ``uh''), supported by transcripts and acoustic feature such as abrupt pitch or intensity variation. \\

\textbf{Positive Attributes} &
Identification of perceptual strengths, such as appropriate emphasis, smooth intonation, or natural fillers, grounded in specific vowels and their acoustic features. \\

\textbf{Potential Issue} &
Identification of perceptual weaknesses, including intonation instability, pacing inconsistency, expressive incorrectness, or unnatural pauses, grounded in specific vowels and their acoustic features.  \\

\bottomrule
\end{tabular}
\label{table:cot-dimensions}
\vspace{-1em}
\end{table*}
